\documentclass[12pt,a4paper]{article}
\usepackage{fontspec}
\setmainfont{Times New Roman}
\setsansfont{Arial}
\setmonofont{Consolas}[Scale=0.85]
\usepackage[left=2.7cm,right=2.3cm,top=2.5cm,bottom=2.5cm]{geometry}
\usepackage{amsmath,amssymb,booktabs,array,tabularx,longtable}
\usepackage{graphicx,float,setspace,enumitem,xcolor,fancyhdr}
\usepackage{caption}
\usepackage{xurl}
\usepackage[colorlinks=true,linkcolor=blue!45!black,citecolor=blue!45!black,urlcolor=blue!45!black]{hyperref}
\hypersetup{pdftitle={Phân tích dữ liệu cảm biến, dự đoán RUL và cảnh báo sớm trên NASA C-MAPSS FD001--FD004},pdfauthor={}}
\renewcommand{\contentsname}{Mục lục}
\renewcommand{\listfigurename}{Danh mục hình}
\renewcommand{\listtablename}{Danh mục bảng}
\renewcommand{\figurename}{Hình}
\renewcommand{\tablename}{Bảng}
\renewcommand{\refname}{Tài liệu tham khảo}
\newcommand{\reportdate}{25/09/2026}
\newcommand{\filepath}[1]{\nolinkurl{#1}}
\newcolumntype{L}[1]{>{\raggedright\arraybackslash}p{#1}}
\onehalfspacing
\setlength{\parindent}{0.7cm}
\setlength{\emergencystretch}{2em}
\setlength{\headheight}{16pt}
\setlist{nosep,leftmargin=*}
\captionsetup{font=small,labelfont=bf}
\pagestyle{fancy}
\fancyhf{}
\fancyhead[L]{\small NASA C-MAPSS FD001--FD004}
\fancyhead[R]{\small Kết quả phân tích -- \reportdate}
\fancyfoot[C]{\thepage}

\begin{document}
\hypersetup{pageanchor=false}
\begin{titlepage}
\centering
{\large TRƯỜNG ĐẠI HỌC CÔNG NGHIỆP TP. HỒ CHÍ MINH\par}
\vspace{1.8cm}
{\Large\bfseries BÁO CÁO CUỐI KỲ\par}
\vspace{0.4cm}
{\large Môn học: Phân tích dữ liệu lớn\par}
\vspace{1.5cm}
{\LARGE\bfseries Phân tích dữ liệu cảm biến, dự đoán RUL và cảnh báo sớm\par}
\vspace{0.6cm}
{\Large NASA C-MAPSS FD001--FD004\par}
\vspace{0.8cm}
{\large Đề tài 8: Phân tích dữ liệu IoT và bảo trì dự đoán\par}
\vspace{1.5cm}
\begin{tabular}{ll}
Học viên thực hiện: & Phan Nguyễn Mai Phương\\
Mã số học viên: & 25002711\\
Ngành: & Khoa học máy tính
\end{tabular}
\vfill
{\normalsize TP. Hồ Chí Minh, tháng 9 năm 2026\par}
\end{titlepage}
\hypersetup{pageanchor=true}
\pagenumbering{roman}
\section*{Tóm tắt}
\addcontentsline{toc}{section}{Tóm tắt}
Báo cáo nghiên cứu một pipeline phân tích dữ liệu cảm biến và ước lượng
Remaining Useful Life (RUL, tuổi thọ hữu ích còn lại) trên bộ dữ liệu mô phỏng
NASA C-MAPSS FD001--FD004, trong đó FD001 là subset chính cho phân tích chi tiết.
Mục tiêu là kiểm tra tính hợp lệ của dữ liệu, xây dựng biểu diễn trailing window,
đánh giá mô hình hồi quy theo đơn vị động cơ và mô tả diễn
biến cảm biến bằng phân tích thống kê cùng PCA--ARIMA. Phương pháp gồm kiểm tra
data contract, ETL qua các lớp Silver/Gold, tạo đặc trưng từ window 30 cycles,
GroupKFold theo engine và đánh giá cuối trên 100 test engines của FD001. Mô hình
HistGradientBoosting đạt mean engine-level MAE bằng 15,3490 cycles, RMSE bằng
20,5413 cycles và NASA score bằng 1.513,2577 trong evaluation protocol được xác định trước. Phân tích
PCA--ARIMA tạo 704 dự báo Health Index trên 10 engine train được chọn trước;
phân tích này được xem là nhánh phụ và không phải mô hình dự báo RUL. Với đầu
ra cảnh báo sớm minh họa, ngưỡng cố định $\mathrm{RUL}\leq 30$ cycles trên OOF
train tạo cảnh báo ở 100/100 engine: 44 trước ngưỡng thật, 7 đúng ngưỡng và 49
sau ngưỡng. Lead time có trung vị 30 cycles và miền 19--55. Cùng OOF đã dùng
để chọn model, nên đây là mô tả nội bộ, không phải đánh giá độc lập về hiệu
quả cảnh báo hoặc lịch bảo trì tối ưu.

Thực nghiệm bổ sung dùng cùng window 10 cycles cho FD001--FD004 để giữ đủ các
test engines. Với model train trên FD001 và áp dụng trực tiếp sang dataset khác,
mean engine MAE lần lượt là 60,3215 cycles trên FD002, 47,3949 trên FD003 và
58,5731 trên FD004. Các giá trị này cao hơn evaluation nội bộ trên từng subset
theo cùng protocol window 10, cho thấy khả năng chuyển miền của model FD001 còn
hạn chế. Metric bổ sung không được so trực tiếp với kết quả chính dùng window
30 cycles.

Kết quả cho thấy pipeline có thể tái lập từ raw FD001 và ma trận bổ sung FD002--FD004
giữ nguyên ranh giới train/test trong quá trình xử lý. Tuy nhiên, FD001 chỉ đại
diện cho một điều kiện vận hành và một dạng suy giảm; kết quả transfer cho thấy
khác biệt giữa các subset trong môi trường mô phỏng. Vì dữ liệu không
cung cấp chi phí bảo trì, downtime hoặc quy đổi cycle sang thời gian lịch, các
kết quả chưa đủ để suy ra chính sách bảo trì hoặc khả năng khái quát cho động cơ
thực tế.

\noindent\textbf{Từ khóa:} NASA C-MAPSS FD001--FD004; RUL; early warning; rolling window;
GroupKFold; ETL; Health Index; PCA--ARIMA.
\clearpage
\begin{spacing}{1.1}
\tableofcontents
\end{spacing}
\clearpage
\input{sections/glossary}
\listoffigures
\listoftables
\clearpage
\pagenumbering{arabic}
\input{sections/problem-data}
\clearpage
\input{sections/theory}
\clearpage
\input{sections/methodology}
\clearpage
\input{sections/statistical-protocol}
\clearpage
\input{sections/current-results}
\clearpage
\section{Thảo luận, giới hạn và hướng phát triển}

\subsection{Diễn giải các kết quả chính}
Kết quả cho thấy việc đặt engine làm đơn vị chia dữ liệu và đánh giá giúp hạn chế việc
đếm các window có phụ thuộc thời gian như những quan sát độc lập. Trong
protocol đánh giá, HistGradientBoosting đạt MAE trung bình theo engine thấp nhất
trên đánh giá chéo và được huấn luyện lại trước khi tạo một dự đoán cuối cho mỗi
engine test. Metric cuối vì vậy phản ánh sai số ở điểm cắt của 100 engine test,
không phải sai số của toàn bộ chuỗi theo thời gian.

Phân tích đầu/cuối vòng đời cho thấy nhiều sensor có thay đổi có hệ thống trong
hai đoạn 30 cycles. Tuy nhiên, độ lớn chênh lệch nằm ở đơn vị gốc của từng
sensor, nên không thể dùng các giá trị này để xếp hạng trực tiếp giữa các
sensor. Wilcoxon signed-rank được diễn giải dưới giả định phân phối chênh lệch
đối xứng quanh 0; không nên rút gọn thành kiểm định median cho mọi phân phối.
Với sensor 6, 53 chênh lệch bằng 0 bị loại khỏi phép xếp hạng theo
\texttt{zero\_method=wilcox}, còn 47 chênh lệch dương tham gia kiểm định.
P-value nhỏ không chứng minh quan hệ nhân quả hay giá trị dự báo ngoài mẫu.

PCA--ARIMA cung cấp một cách kiểm tra độc lập về diễn biến của Health Index.
Tỷ lệ bao phủ của khoảng dự đoán 95\% thay đổi đáng kể giữa các engine, từ 21,43\%
đến 97,03\%. ADF trên các đoạn huấn luyện không bác bỏ giả thuyết gốc về nghiệm
đơn vị; do đó, các chuỗi này chưa cung cấp bằng chứng đủ mạnh về tính dừng theo
kiểm định ADF. Sau khi điều chỉnh bậc tự do theo bậc ARIMA, Ljung--Box vẫn cho
thấy p-value dưới 0,05 ở engine 69 và 70. Dự báo Health Index vì vậy chỉ được
giữ như một phân tích phụ, không được dùng để thay thế mô hình RUL.

So sánh với baseline trên cùng prefix và suffix cho thấy ARIMA có MAE trung bình
0,4110, thấp hơn \texttt{drift} (0,7469) và \texttt{last\_value} (0,9478); RMSE tương ứng là
0,5027, 0,9128 và 1,1553. ARIMA chỉ thấp hơn \texttt{drift} ở 7/10 engine và
\texttt{last\_value} ở 8/10 engine, nên kết quả này phù hợp với một lợi thế mô tả trong
subset được chọn hơn là bằng chứng về ưu thế tổng quát. Baseline không có
khoảng dự đoán, vì vậy không được so tỷ lệ bao phủ với ARIMA.

Cảnh báo sớm minh họa trên OOF cho thấy ngưỡng 30 cycles có thể chuyển
dự đoán thành một sự kiện theo engine thay vì đếm từng cycle như các quan sát
độc lập. Tuy nhiên, 49/100 cảnh báo đầu tiên đến sau mốc RUL thật bằng 30;
lead time thấp nhất là 19 cycles. Không có engine hoàn toàn thiếu cảnh báo
không đồng nghĩa mọi cảnh báo đều đủ sớm. Các nhóm trước/đúng/sau ngưỡng chỉ
mô tả thời điểm, chưa phản ánh chi phí hoặc giá trị sử dụng cho bảo trì. Không có engine
âm tính theo định nghĩa sự kiện toàn vòng đời, nên tỷ lệ cảnh báo sai cấp engine
không xác định. Vẫn có thể định nghĩa sai cảnh báo theo từng vùng thời gian,
nhưng không được thay thế âm thầm đại lượng cần ước lượng hiện tại bằng phép đo đó.

Kết quả warning còn chịu ảnh hưởng của model selection: cùng tập OOF đã dùng
để chọn HistGradientBoosting rồi mô tả event. Dự đoán mỗi engine không dùng
engine đó trong fit, nhưng toàn bộ phân tích chưa độc lập với lựa chọn model.
Muốn ước lượng hiệu quả của cả quy trình lựa chọn cần một đánh giá vòng ngoài
riêng; 100\% từng phát cảnh báo không thay thế được đánh giá này.

\subsection{Giới hạn phương pháp}
Thứ nhất, FD001 chỉ gồm một điều kiện vận hành và một dạng suy giảm mô phỏng.
Ma trận bổ sung với window 10 cycles cho thấy MAE transfer từ FD001 tăng lên
60,3215 cycles trên FD002, 47,3949 trên FD003 và 58,5731 trên FD004, trong khi
evaluation nội bộ của các subset nằm khoảng 21,9258--28,3255 cycles. Kết quả này
cho thấy khác biệt phân bố giữa các subset có ảnh hưởng đáng kể trong protocol đang xét; nó không
cho phép suy ra một metric chung cho mọi điều kiện. Thứ hai, test evaluation chỉ có một
điểm cắt cho mỗi engine; cách đánh giá này không mô tả sai số khi mô hình được
gọi lặp lại ở nhiều cycle. Thứ ba, runtime benchmark hiện chỉ đo Pandas trên
dữ liệu nhân bản trong bộ nhớ, với hai lần đo mỗi scale. PySpark mới được
kiểm tra đối chiếu: kết quả tương đương trên input đã xét không phải bằng chứng
về tốc độ hoặc khả năng mở rộng của hệ thống phân tán.
Cuối cùng, dữ liệu không cung cấp chi phí bảo trì, downtime hoặc quy đổi cycle
sang thời gian lịch; vì vậy báo cáo chưa thể chuyển metric RUL thành quyết định
bảo trì.

\subsection{Hướng nghiên cứu tiếp theo}
Khả năng khái quát có thể được khảo sát thêm bằng cách thay đổi độ dài window,
chia dữ liệu theo thời gian trên nhiều engine và dùng đánh giá vòng ngoài cho
quy trình chọn model. Với Health Index, cần huấn luyện transformer ngoài các engine
được dự báo và báo cáo đồng thời tỷ lệ phát hiện lỗi, tỷ lệ bao phủ và khoảng thời gian dự báo.

Kết quả FD002--FD004 hiện đã cung cấp một phép kiểm tra chuyển sang subset khác.
Các bước tiếp theo có thể đánh giá huấn luyện lại trên từng subset, chuẩn hóa theo
operating condition hoặc điều chỉnh theo miền dữ liệu. Mô hình LSTM/Transformer cần một
giả thuyết về biểu diễn thời gian và so sánh cùng protocol với baseline hiện
tại. Khi thiếu cost, downtime và intervention, chưa thể chuyển kết quả RUL
thành chính sách bảo trì.

\clearpage
\section{Kết luận}

Báo cáo trả lời năm câu hỏi nghiên cứu chính trên NASA C-MAPSS FD001 và một câu
hỏi bổ sung về chuyển miền sang FD002--FD004.
Thứ nhất, FD001 có 20.631 train observations và 13.096 test observations đáp ứng các
kiểm tra schema, khóa và tính liên tục của cycle; điều này xác nhận khả năng
xử lý dữ liệu theo contract, không xác nhận tính đại diện cho động cơ thực tế.
Thứ hai, phân tích đầu/cuối 30 cycles ghi nhận 15 sensor bác bỏ giả thuyết
không của Wilcoxon sau hiệu chỉnh BH trong cấu hình đã chọn; kết quả cần đọc
cùng giả định signed-rank, độ lớn chênh lệch và số zero, không dùng làm
feature selection ngoài training fold.

Thứ ba, Pandas và PySpark khớp ba boundary Silver/Gold đã kiểm tra, với sai
khác tuyệt đối lớn nhất bằng 0. Runtime chỉ được đo cho Pandas; báo cáo chưa
có bằng chứng so sánh tốc độ hai engine hoặc hiệu năng cluster. Thứ tư,
HistGradientBoosting có mean engine MAE thấp nhất trong ba ứng viên trên
GroupKFold. Sau refit, model đạt MAE 15,3490 cycles và RMSE 20,5413 cycles tại
điểm cắt của 100 test engines; bootstrap CI 95\% cho MAE là
[12,7286; 18,0866]. Metric tại điểm cắt không mô tả toàn bộ trajectory.

Thứ năm, ngưỡng cảnh báo minh họa 30 cycles tạo ít nhất một event trên mỗi
engine OOF train, nhưng chỉ 44 engine cảnh báo trước ngưỡng và 7 đúng ngưỡng;
49 engine cảnh báo sau ngưỡng. Lead time có trung vị 30 cycles, khoảng tứ
phân vị 26--35 và miền 19--55. Vì OOF cũng dùng để chọn model, đây là đánh giá
nội bộ, chưa chứng minh chất lượng cảnh báo độc lập hay giá trị sử dụng cho bảo trì.
Ví dụ engine 1 cảnh báo tại cycle 162, trước failure proxy cycle 192 đúng
30 cycles; điều đó chưa quy định cần can thiệp khi nào.

Nhánh phụ PCA--ARIMA có MAE trung bình 0,4110, thấp hơn drift 0,7469 và
last-value 0,9478 trên cùng 10 chuỗi giữ lại để kiểm tra. Tuy nhiên, ARIMA chỉ có
MAE thấp hơn drift ở 7/10 engine, và tỷ lệ bao phủ thay đổi từ 21,43\% đến 97,03\%.
Kết quả này hỗ trợ phân tích Health Index, không thay thế dự đoán RUL.

Đóng góp của báo cáo là một pipeline có thể tái lập, các phép so sánh được mô tả
theo cùng evaluation protocol và các giới hạn được công bố. Ma trận bổ sung với
window 10 cycles cho thấy model train trên FD001 có MAE 60,3215 trên FD002,
47,3949 trên FD003 và 58,5731 trên FD004, cao hơn evaluation nội bộ của các
subset. Vì vậy, kết luận chính được giữ ở phạm vi FD001; các subset khác cho
thấy cần huấn luyện lại hoặc xử lý khác biệt giữa các subset trước khi áp dụng.
Phạm vi dữ liệu hiện tại chưa cho phép kết luận về lịch bảo trì tối ưu hoặc hiệu
quả triển khai thực tế.

\clearpage
\appendix
\input{sections/appendix}
\clearpage
\phantomsection
\addcontentsline{toc}{section}{Tài liệu tham khảo}
\bibliographystyle{plainurl}
\bibliography{../references/bibliography}
\end{document}
